\section{Erd\H{o}s Problem \#308: the first missing integer Egyptian fraction}
\noindent Reconstruction of the asymptotic result and its indexing, 23 September 2026.

\subsection*{1) Statement and separate extremal quantities}
Let $\mathcal R_N$ be the positive integers representable as $\sum_{d\in S}1/d$ with $S\subseteq\{1,\ldots,N\}$. Put $H_N=\sum_{d=1}^N1/d$ and
\[
 f(N)=\min(\mathbb N\setminus\mathcal R_N),\qquad
 h(N)=\max\{h:\{1,\ldots,h\}\subseteq\mathcal R_N\}.
\]
Thus $f(N)=h(N)+1$. The question asks about this first missing integer and whether $\mathcal R_N$ is always an initial interval. We establish the complete eventual-interval conclusion using Croot's main theorem, with its external status explicit. We do not establish the interval assertion for every finite $N$.

\subsection*{2) The precise external theorem and an indexing correction}
In \emph{On some questions of Erd\H{o}s and Graham about Egyptian fractions}, Croot defines his extremal quantity as the largest integer up to which \emph{every} positive integer is representable. It is $h(N)$ above, not the first missing integer. His Main Theorem, on page 1 of the author's manuscript, gives, with $E_N=(\log\log N)^2/\log N$,
\[
 \left\lfloor H_N-(9/2+o(1))E_N\right\rfloor
 \le h(N)\le
 \left\lfloor H_N-(1/2+o(1))E_N\right\rfloor.\tag{1}
\]
The two error terms need not be the same. The analytic construction and obstruction proving (1) are external inputs. Accordingly the bounds for the quantity actually called $f(N)$ here are
\[
 \left\lfloor H_N-(9/2+o(1))E_N\right\rfloor+1
 \le f(N)\le
 \left\lfloor H_N-(1/2+o(1))E_N\right\rfloor+1.\tag{2}
\]
Keeping this $+1$ is necessary. Applying (1) directly to the first missing integer changes its meaning.

\subsection*{3) Why there are eventually no internal gaps}
Every represented integer is at most $H_N$, so $\mathcal R_N\subseteq\{1,\ldots,m_N\}$, where $m_N=\lfloor H_N\rfloor$. Since $E_N\to0$, the subtracted term in the lower bound of (1) lies between zero and one for all sufficiently large $N$. Thus
\[
 h(N)\ge m_N-1.
\]
Every integer through $m_N-1$ is therefore represented; only $m_N$ itself can be in doubt. Consequently
\[
 \mathcal R_N=\{1,\ldots,m_N-1\}
 \quad\text{or}\quad
 \mathcal R_N=\{1,\ldots,m_N\},\tag{3}
\]
and $f(N)$ is respectively $m_N$ or $m_N+1$. This conclusion requires no downward-closure assumption about arbitrary Egyptian fraction representations.

Writing $H_N=m_N+\theta_N$ with $0\le\theta_N<1$ clarifies the rounding. The upper bound in (1) excludes $m_N$ whenever $\theta_N<(1/2+o(1))E_N$; its lower bound includes $m_N$ whenever $\theta_N\ge(9/2+o(1))E_N$. The narrow intermediate range is not decided by these two estimates alone. In either case $f(N)=\log N+O(1)$ by the elementary estimate $H_N=\log N+\gamma+O(1/N)$.

\subsection*{4) Exact small examples and limits of this reconstruction}
For $1\le N\le5$, the fractions with denominators at least two cannot sum to an integer. For $N=5$, a putative sum one from $\{1/2,1/3,1/4,1/5\}$, multiplied by 60, would read $30a+20b+15c+12d=60$, with $a,b,c,d\in\{0,1\}$. Reduction modulo 5 gives $d=0$, modulo 3 then gives $b=0$, and $2a+c=4$ is impossible. A sum two is excluded by their total $77/60<2$. Smaller $N$ are subsets of this case. Hence $\mathcal R_N=\{1\}$ and $f(N)=2$ for $N\le5$. At $N=6$, the identity $1/2+1/3+1/6=1$ supplies the integer two when combined with denominator one; $H_6<3$ excludes any larger integer. Thus $\mathcal R_6=\{1,2\}$ and $f(6)=3$.

\textbf{UNRESOLVED in the literal all-$N$ formulation in this attempt.} The eventual interval theorem, correct first-missing indexing, and asymptotic estimates are established relative to (1). No proof eliminating every possible small exceptional $N$, exact decision rule for the transition range, or independent proof of Croot's analytic input is supplied.

Sources: \url{https://www.erdosproblems.com/308}, checked 23 September 2026; E. S. Croot III, Mathematika 46 (1999), 359--372, author's primary manuscript \url{https://ecroot.math.gatech.edu/erdos_graham.pdf}, Main Theorem on page 1. The original definition and formula were checked visually because the PDF text encoding is defective.

The estimate measures coverage of the question, including the remaining all-$N$ qualification, not a correctness probability.
\subsection*{5) COMPLETION ESTIMATE}
\textbf{COMPLETION: 80\%}
